\section{Dependent types: further details}
\applabel{dependent types}

We begin by establishing the correctness of the action over the
Grothendieck construction $\coprod_{\carrier\Ctx}\A$, as we will use
its properties pervasively, e.g., when extending contexts.

\begin{proof}[of \propositionref{Grothendieck action}]%
  Take a $\Ctx$-indexed action $\A$ and define an action on
  $\coprod_{\carrier\Ctx} \carrier\A$ by:
  \[
  g \act \pair{\ctx : \carrier\Ctx}{a : \carrier\A_{\ctx}}
  \definedby \pair{g \act \ctx}{g \iact{\ctx} a : \carrier\B (g \act \ctx)}
  \]
  Since $\projection_1(g \act\pair\ctx a) = g \act \ctx = \projection_1\pair\ctx a$,
  the projection is equivariant
  $\projection_1 : \coprod_{\Ctx}\A \equito \Ctx$.

  Conversely, given an action $\coprod_{\Ctx}\A$ as in the assumption,
  define $g \iact{\ctx} a \definedby \projection_2(g \act
  \pair{\ctx}{a})$. This definition is a well-typed dependent function $\ctx : \Ctx \types (\iact{}) : (g : \carrier\G) \to A\ctx \to A (g \act \ctx)$:
  \[
  g \iact\ctx a \in A\parent{\projection_1(g \act \pair \ctx a)} =
  A\parent{\g \act \ctx}
  \]
  It also satisfies the two conditions for an indexed action. The unit axiom:
  \[
  \unit \iact \ctx a = \projection_2 (\unit \act \pair \ctx a) = \projection_2\pair\ctx a = a
  \]
  The multiplication axiom:
  \begin{align*}
  h \iact{g \act \ctx} (g \iact \ctx a)
  &=
  \projection_2\parent{h \act
    \pair{g\act \ctx}
         {\projection_2\parent{g \act \pair \ctx a}}}
  \\&
  \explain={$\projection_1$ equivariant}
  \projection_2\parent{h \act
    \pair{\projection_1\parent{g\act \pair\ctx a}}
         {\projection_2\parent{g \act \pair \ctx a}}}
  \\&
  \explain={extensionality of pairs}
  \projection_2\parent{h \act \parent{g \act \pair \ctx a}}
  =
  \projection_2\parent{(h \mul g) \act \pair \ctx a}
  =
  (h \mul g) \iact\ctx a
  \end{align*}
  and so we have an indexed action. The two constructions are inverse
  to each other.
\end{proof}

The following example lets us work more formally with concrete
examples:

\begin{example}\examplelabel{equivariant Grothendieck}
The Grothendieck construction extends equivariantly to equivariant
terms and maps.  It sends equivariant terms $\Ctx \types t : \A$ to
equivariant sections \((\fun \ctx \pair{\ctx}{t_{\ctx}}) : \Ctx
\equito \coprod_{\Ctx}\A \). It sends a dependent equivariant function
$\Ctx \types f : \DFun\x\A \B$ to the following equivariant map
$\parent{\fun{\pair\ctx a}\pair{\ctx, \For x a}{f_{\ctx}a}} :
\coprod_{\Ctx}\A \equito \coprod_{\Ctx, \x : \A}\B$.
\end{example}

Next, we establish the indexed action axioms for the family of
dependent functions, and the indexed equivariance of related
operations.

\begin{lemma}\lemmalabel{dependent function action}
  The indexed action structure over the family of dependent functions
  is an action, evaluation is equivariant, and currying and uncurrying
  preserve equivariance.
\end{lemma}
\begin{proof}
  Indeed, the action axioms hold by calculating with the correct
  indices as in the ordinary case. Similarly for evaluation:
  \[
  \eval(g\iact{\ctx} f,\x)_{g\act \ctx, \For \x g\iact\ctx a}
  =
  g \iact{} f(\inv g \iact{} g \iact{} a)
  =
  g \iact{} fa
  =
  g\iact{}\eval(f, \x)_{\ctx, \For\x a}
  \]
  The bijective correspondence is given as in the set-theoretic
  case. If $t$ is equivariant, then:
  \begin{align*}
  (\fun{\x:\A}t)_{g \act \ctx}a &=
  g \iact{} t_{g\act \ctx, \For\x a}
  =
  g \iact{} t_{g\act \ctx, \For\x (g \act (\inv g \act a))}
  \\&\explain={equivariance}
  g \iact{} g\iact{} t_{\ctx, \For\x (\inv g \act a)}
  =
  (g \iact{} (\fun{\x : \A}t)_{\ctx})a
  \end{align*}
  Therefore $(\fun{\x:\A}t)_{g \act \ctx} = (g \iact{} (\fun{\x : \A}t)_{\ctx})$.

  Conversely, if $\fun{(\x:\A)}t$ is equivariant, then since
  $\eval\pair{-}\x$ is equivariant, $t$ is.
\end{proof}



We now turn to prove the indexed counterpart to the \WLOG{}
Representation \theoremref{wlog-construction}:

\begin{proof}[of \theoremref{indexed-wlog-construction}]%
  We can prove this theorem directly, but it may be more interesting
  to reduce it to the ordinary representation theorem. Applying the
  Grothendieck construction to the situation, we have:
  \begin{itemize}
  \item
    $\G$-actions $\coprod_{\Ctx}\A$ and $\coprod_{\Ctx, \x : \A}\B$;
  \item canoniser $\canoniser' \definedby \fun{\pair\ctx a}\canoniser_{\ctx}a :
    \carrier{\coprod_{\Ctx}\A} \to \carrier\G$;
  \item below we show that $\Canon{\canoniser'} = \set{\pair\ctx a \suchthat
    a \in \Canon{\canoniser, \ctx}} \defines \coprod_{\Ctx}\Canon{\canoniser}$;
  \item core function $f' \definedby
    \parent{\fun{\pair\ctx a}\pair{(\ctx, \For\x a)}{f_{\ctx}a}}
    : \Canon{\canoniser'} \to \carrier{\coprod_{\Ctx,\x:\A}\B}$;
  \item we have that $\wlg(\canoniser', f'; \pair\ctx a) =
    \pair{(\ctx, \For\x a)}{\wlg(\canoniser, f)_{\ctx}a}$;
  \item the core function $f'$ is equivariant iff $f$ is
    (by \exampleref{equivariant Grothendieck});
  \item below we show that $\wlg(\canoniser', f')$ is equivariant iff
    $\wlg(\canoniser, f)$ is equivariant;
  \item below we show that the function term $\carrier\Ctx \types g : \DFun\x{\carrier\A}{\carrier\B}$
    extends $f$ along $\Ctx \types \Canon{\canoniser} \subset \A$ iff
    the function $\fun {\pair\ctx a}\pair-{(\ctx, \For\x a)}{g_{\ctx}a}
    : \coprod_{\carrier\Ctx}{\carrier\A}
    \to \coprod_{\carrier\Ctx, \x : \carrier\A}\carrier\B$ extends
    $f'$ along $\Canon{\canoniser'}\subset
    \coprod_{\carrier\Ctx}\carrier\A$.
  \end{itemize}
  The theorem follows by unpacking and re-packing definitions in both directions. Let's discharge the outstanding proof obligations.

  Let $\pair\ctx a$ be canonical for $\canoniser'$, so $\pair\ctx a =
  (\canoniser_{\ctx}a') \act\pair{\ctx'}{a'}$. Then $\ctx' =
  \inv*{\canoniser a'}\act \ctx$ and $a' \in \A_{\inv*{\canoniser
      a'}\act \ctx}$, so $a \in (\Canon\canoniser)_{\ctx}$, and so
  $\pair\ctx a \in \coprod_{\Ctx}\Canon{\canoniser}$. These
  implications are invertible, proving the converse.

  Next, we show that the Grothendieck construction preserves the \WLOG
  extension:
  \begin{align*}
  \wlg(\canoniser', f'; \pair\ctx a) &=
  \inv*{\canoniser'\pair\ctx a} \act
    f'(\parent{\canoniser'\pair\ctx a}\act \pair\ctx a)
  \\&=
  \inv*{\canoniser_{\ctx} a} \act
    f'((\canoniser_{\ctx} a)\act \pair\ctx a)
  \\&=
  \inv*{\canoniser_{\ctx} a} \act
    f'\pair{(\canoniser_{\ctx} a)\act\ctx}{(\canoniser_{\ctx} a) \iact{} a}
  \\&=
  \inv*{\canoniser_{\ctx} a} \act
  \pair{\parent{(\canoniser_{\ctx} a)\act\ctx,
                  \For\x((\canoniser_{\ctx} a) \iact{} a)}}
       {f_{(\canoniser_{\ctx} a)\act\ctx}{(\canoniser_{\ctx} a) \iact{} a}}
  \\&=
  \pair{\parent{\ctx,
                  \For\x a}}
       {\inv*{\canoniser_{\ctx} a} \act
         f_{(\canoniser_{\ctx} a)\act\ctx}{(\canoniser_{\ctx} a) \iact{} a}}
  \\&=
  \pair{\parent{\ctx,
                  \For\x a}}
       {\wlg(\canoniser, f)_{\ctx}a}
  \end{align*}
  In addition, by \exampleref{equivariant Grothendieck}, one \WLOG
  extension is equivariant iff the other is equivariant.

  The final bullet point follows by functoriality of the Grothendieck
  construction.
\end{proof}

We now establish the bread-and-butter properties facilitating smooth
dependently-typed symmetric programming and symmetric propositional
reasoning. We start with dependent-function abstraction:

\begin{lemma}[abstraction]
  Function abstraction of an equivariant term $\Ctx, \x : \A \types t
  : \B$ is both:
  \begin{itemize}
  \item an equivariant term $\Ctx \types \parent{\fun{(\x:\A)}t} : \DFun\x\A\B$.
  \item an equivariant function $\Ctx \types \parent{\fun{(\x:\A)}t} : \EFun\x\A\B$.
  \end{itemize}
\end{lemma}
\begin{proof}
  Calculate, ignoring indices (as those type-check):
  \begin{align*}
    \parent{g\act(\fun{\x:\A}t)}_{\ctx}\ma
    &\explain{=}{action on\\dep. functions}
  g \act ((\fun{\x:\A}t)(\inv g \act \ma))
  \explain={abstraction def}
  g \act t_{\ctx, \For\x(\inv g \act \ma)}
  \explain={$t$ equivariant}
  t_{g \act \ctx, \For\x(g \act \inv g \act \ma)}
  \\&=
  t_{g \act \ctx, \For\x \ma}
  =
  (\fun{\x : \A}t)_{g\act\ctx} \ma
  \end{align*}
  showing abstraction is an equivariant term. It is also an
  equivariant function, for:
  \[
  (\fun{\x:\A}t)_{g\act\ctx}(g \act \ma)
  \explain={abstraction def}
  t_{g\act\ctx, \For \x{g\act\ma}}
  \explain={$t$ equivariant}
  g \act t_{\ctx, \For\x\ma}
  \]
  as we wanted.
\end{proof}

Substitution using an equivariant substitution preserves equivariance:

\begin{lemma}[reindexing and substitution]
Reindexing by an equivariant function $\msubst : \Ctx
\to \Dtx$ sends equivariant terms $\Dtx \types s : \B$ to equivariant
terms $\Ctx \types \subst s\msubst : \subst \B\msubst$, and moreover
every equivariant term $\Ctx \types t : \subst\B\msubst$ induces an
equivariant substitution $(\msubst, \For\y t) : \Ctx \equito (\Dtx, \y :
\B)$.
\end{lemma}
\begin{proof}
  Calculate:
  \[
    g\act \subst s\msubst_{\ctx}
    \explain={reindexing def}
    g\act s_{\msubst\ctx}
    \explain*={$s$ equivariant}
    s_{g\act \msubst\ctx}
    \explain={$\msubst$ equivariant}
    s_{\msubst(g\act \ctx)}
    =
    \subst s\msubst_{g\act \ctx}
  \]
  and so $\subst s\msubst$ is equivariant. Next, calculate:
  \begin{align*}
  (\msubst, \For\y t)(g \act \ctx)
  &\explain={substitution def}
  (\msubst(g\act\ctx), \For\y t_{g \act \ctx})
  \explain={$t$ equivariant}
  (\msubst(g\act\ctx), \For\y g\act t_{\ctx})
  \explain={action on context extension}
  g\act (\msubst\ctx, \For\y t_{\ctx})
  \\&=
  g\act((\msubst, \For\y t)\ctx)
  \end{align*}
  and so $(\msubst, \For\y t) : \Ctx \equito (\Dtx, \y : \B)$.
\end{proof}

Dependent pairing also preserves equivariant terms:

\begin{lemma}[dependent pairing]
  Two equivariant terms $\Ctx \types t : \A$ and $\Ctx \types s :
\subst\B{\For\x t}$ induce an equivariant term $\Ctx \types \dpair ts
: \DPair\x\A\B$.
\end{lemma}
\begin{proof}
  Calculate:
  \[
  (g \act (\dpair ts))_{\ctx}
  =
  (\dpair{g \act t_{\ctx}}{g \act s_{\ctx, \For\x t_{\ctx}}})
  \explain={$t$, $s$ equivariant}
  (\dpair{t_{g \act \ctx}}{s_{g\act\ctx, \For\x (g \act t_{\ctx})}})
  \]
  as we wanted.
\end{proof}

We now turn to propositions:

\begin{lemma}[subfamilies]
Each equivariant proposition $\Ctx, \x : \A \types \phi : \Prop$
induces an indexed action on the subfamily $\Ctx \types \set{\x : \A
  \suchthat \phi}$ that lifts $\A$ along the subset inclusions.
Equivariant terms $\Ctx \types t : \A$ satisfying $\phi$ lift to
equivariant terms inhabiting the associated subspace i.e., $\Ctx
\types \lift t : \set{\x : \A \suchthat \phi}$.
\end{lemma}
\begin{proof}
  Equivariance means that if $a \in \carrier\A_{\ctx}$ satisfies
  $\phi$, then so does $g \act a$. Therefore $g \iact\ctx a \in
  \set{\x : \A \suchthat \phi}$, and the action lifts.  For
  equivariance of the lifted term, use injectivity of the coercion
  function $\Ctx \types \coerce \definedby \fun \ma \ma : \set{\x : \A
    \suchthat \phi} \equito \A$:
  \[
  \coerce_{g \act \ctx} (g \act (\lift t)_{\ctx}) \explain={$\coerce$ equivariance}
  g \act \coerce_{\ctx} \lift t_{\ctx} \explain={coercion}
  g \act t_{\ctx} \explain={$t$ equivariant}
  t_{g\act\ctx} =
  \coerce_{g\act \ctx} (\lift t)_{g\act\ctx}
  \]
  and so $g \act (\lift t)_{\ctx} = (\lift t)_{g\act\ctx}$.
\end{proof}

We use these facts implicitly when conducting reasoning:

\begin{lemma}
The logical connectives are equivariant $(\land), (\lor) : \Prop^2 \equito \Prop$,
$(\lnot) : \Prop \equito \Prop$, as well as the quantifiers $\forall,\exists :
(\A \to\Prop) \equito \Prop$.
\end{lemma}
\begin{proof}
  The logical connectives operate between trivial actions hence
  equivariant.  For the quantifiers, assume $\forall t = \Tru$, and
  so, for each $\ma \in \carrier\A_{\ctx}$, we have $t(\inv g \act
  \ma) = \Tru$, i.e.:
  \[
  \forall (g \act t) = \forall \x . t(\inv g \act
  \x) = \Tru
  \] Conversely, if $\forall (g \act t)$, then by
  instantiating with $g \act \ma$ we have that $t\ma = t(\inv g \act g
  \act \ma) = \Tru$. Therefore $\forall t = \forall g \act t$. A
  similar argument holds for $\exists$, or by use of de Morgan.
\end{proof}

We conclude our treatment of propositions by explicating how symmetric
reasoning \WLOG{} often manifests---via a property invariant under the
symmetry of interest:

\begin{proof}[of \lemmaref{core functions to propositions invariance}]%
  Take canonical $a$ and $g \iact\ctx a$. By type-checking, $g \iact{}
  f_{\ctx} a, f_{\ctx}(g \iact{} a)$ prove the same proposition,
  namely $\phi_{g\act\ctx, \For\x(g \iact{} a)}$.  Therefore, $g
  \iact{} f_{\ctx} a = f_{\ctx}(g \iact{} a)$.
\end{proof}

In some concrete cases, working with the entire group of permutations
incurs a combinatorial explosion. Group theorists often generate the
group from a smaller set of symmetries, and then only work with those.
%
As we briefly mentioned in the main text,
a group $\G$ is \emph{generated} by a subset $S \subset
\carrier\G$ when $\G$ is the smallest subgroup of $\G$ containing $S$, equivalently,
when every $\carrier\G$-element is the multiplication of
$\carrier\G$-elements or their inverses.
For example, the symmetry group $\Sym 3$ is generated by
the two transpositions $(1 \, 2)$ and $(2 \, 3)$.
We have an indexed variation of a well-known lemma:
\begin{lemma}\lemmalabel{equivariance by generators}
  Let $\Ctx \types \A$ and $\Ctx, \x : \A \types \B$ be indexed
  $\G$-actions.  If $S$ generates $\G$ and $\carrier\Ctx \types f :
  \DFun\x{\carrier\A}{\carrier\B}$ is equivariant w.r.t.~symmetries in
  $S$, then $f$ is equivariant $\Ctx \types f : \EFun\x{\A}{\B}$. When
  $\B = \Prop$ and $\inv{\bracks{S}} \subset S$,
  it suffices to show lax/oplax equivariance, i.e.:
  \[
  \parent{\forall g \in S, \forall \x. f(g\act \x) \implies f \x}
  \lor
  \parent{\forall g \in S, \forall \x. f(g\act \x) \impliedby f \x}
  \]
\end{lemma}
\begin{proof}[of \lemmaref{equivariance by generators}]
  Let $H \subset \carrier\G$ be the set of symmetries that $f$
  respects. Then $S \subset H$. Since $f (e \act a) = f a = e \act (f a)$, we have:
  $e \in H$. Take $h \in H$ we have:
  \[
  f(h \act a) = h \act (fa) \implies f (\inv h \act a) =
  \inv h \act f(h \act (\inv h \act a)) = \inv h \act (f a)
  \]
  so $\inv h \in H$.
  Take $g,h \in H$ and then:
  \[
  f ((g \mul h) \act a) = f (g \act (h \act a)) = g \act (f (h \act
  a)) = g \act (h \act (f a)) = (g \mul h)\act (f a)
  \]
  so $g \mul h \in H$.
  Thus $H \subset \carrier\G$ is a subgroup containing $S$, and so $H
  = \carrier\G$.

  When $\B = \Prop$ and $\inv{\bracks{S}} \subset S$, by substituting
  $\For\x{\inv g \act a}$ we have:
  \[
  f x \iff f (\inv g \act (g \act x))
  \explain\implies{$\inv{\bracks{S}} \subset S$}
  f (g \act x)
  \]
  For the left disjunct, $f$ is in fact equivariant for
  symmetries in $S$.  Similarly for the right disjunct.
\end{proof}

We now turn to study equivariant canonisers. It turns out that the
indexed setting is fruitful for this purpose. We first show that the
indexed action of \emph{distinguishable symmetries} $\G/\A$ from
\definitionref{indistinguishability} is well-defined:

\begin{lemma}\lemmalabel{distinguishable symmetry action well-defined}
  Let $\A$ be a $\G$-action. The action over distinguishable symmetries $\G/\A$
  is congruent w.r.t.~indistinguishability, hence is well-defined.
\end{lemma}
\begin{proof}
  Assume $s_1 \stabeq a s_2$. Calculate:
  \[
  (s_1 \mul \inv g)\act (g \act a) = s_1 \act a = s_2 \act a
  = (s_2 \mul \inv g)\act (g \act a)
  \]
  and so $s_1\mul \inv g \stabeq{g \act a} s_2\mul \inv g$, and the
  action is well-defined.
\end{proof}

In \subsecref{equivariant canonisers} we mentioned in passing the
stabiliser family alongside the indexed actions of the orbits and the
indistinguishable symmetries. The stabilisers also form an indexed
action:
\begin{example}
  The \emph{stabiliser} family has all the symmetries fixing an index
  $a \in \carrier\A$ as the fibre over it: $a : \carrier\A \types
  \stabilise[\G] a \definedby \set{s \in \carrier\G \suchthat s \act a
    = a}$.
  The stabiliser has an $\A$-indexed action:
  \[
  g \iact a{s} \definedby s \mul \inv g \in \stabilise{g \act a}\quad
  \text{since}\quad s \mul \inv g \act (g \act a) = s \act a = a
  \tag*{\qedhere}
  \]
\end{example}

\goodbreak

We now prove our reformulation of the famous Orbit-Stabiliser
\theoremref{orbit stabiliser}:
\begin{proof}[of \theoremref{orbit stabiliser}]%
  The bijectivity of $\phi$ is the well-known theorem. We show equivariance:
  \begin{align*}
  \phi_{g \act x}(g \iact x [h]_{\stabeq x})
  &\explain={action on $\G / \A$} \phi_{g \act x}([h \mul \inv g]_{\stabeq {g \act x}})
  = (h \mul \inv g) \act (g \act x)
  = h \act x
  \explain={action on $\OrbitOf{x}$} g \iact x (h \act x)
  \\&= g \iact x \phi_{x}[h]_{\stabeq x}
  \end{align*}
  and we conclude as inverses of equivariant functions are equivariant.
\end{proof}

Next, we characterise equivariant canonisers as those for which
stable and equivariant core functions coincide. First, an observation:

\begin{lemma}
  \lemmalabel{equivariant canoniser weakly immediate}
  Let $\canoniser : \carrier\A \to \carrier\G$ be an equivariant
  canoniser for a $\G$-action $\A$. For every canonical element $a \in
  \Canon\canoniser$, we have
  \(
  \canoniser a \stabeq a \unit
  \).
\end{lemma}
\begin{proof}
  Since $a$ is canonical, there is some $b \in \carrier\A$ such that
  $a = (\canoniser b)\act b$. Since $\canoniser$ is equivariant, we have:
  \[
  \canoniser a = \canoniser((\canoniser b) \act b)
  \explain*{\stabeq{a}}{equivariance\\$(\canoniser b)\act b = a$}
  (\canoniser b) \mul \inv*{\canoniser b}
  = \unit
  \]
  as we wanted.
\end{proof}

Now, for the proof:

\begin{proof}[\theoremref{all stable core functions are equivariant}]%
  First, assume $\canoniser : (a : \A) \equito (\G/\A)a$ is
  equivariant, and consider any stable core function $f :
  \Canon\canoniser \to \B$. We show that $f$ is equivariant.

  Consider any canonical $g \act a, a \in \Canon{\canoniser}$. Since
  $\canoniser$ is equivariant, by \lemmaref{equivariant canoniser
    weakly immediate}, we have $\canoniser a \stabeq a \unit$ and so
  $\canoniser a \in \stabilise a$.  By equivariance, we have:
  \[
  g \act[\G/\A] [\canoniser a]_{\stabeq a}
  =
  [\canoniser(g \act a)]_{\stabeq{g \act a}}
  \explain={\lemmaref{equivariant canoniser weakly immediate} for $\g \act a$}
  [\unit]_{\stabeq{\g \act a}}
  \implies
  \bracks{\canoniser a}_{\stabeq a}
  = \inv g \act \bracks\unit_{\stabeq{g \act a}}
  = \bracks g_{\stabeq a}
  \]
  and so $g \stabeq a \canoniser a \in \stabilise a$ and so $g \in
  \stabilise a \subset \stabilise {f a}$. Therefore:
  \[
  f(g \act a)
  \explain={$g \in \stabilise a$}
  f a
  \explain*={$g \in \stabilise{f a}$}
  g \act f a
  \]
  and so $f$ is equivariant.

  Conversely, assume that every stable core function is equivariant. We need
  to show the canoniser is equivariant.
  The proof adapts the argument in \theoremref{all core functions are partially equivariant}. Since we now have the indexed action
  vocabulary, the proof simplifies somewhat.

  Take the $\A$-indexed action $\B \definedby \G/\A$, and the following
  core function:
  \[
  f : \DFun\x{\Canon{\canoniser}} (\G/\A)\x
  \qquad
  f a \definedby [\unit]_{\stabeq a}
  \]
  It is stable. Indeed, take any canonical $a \in \Canon{\canoniser}$
  and a stabilising symmetry $g \in \stabilise a$. Since $\stabilise
  a$ is a sub-group, $\inv g \in \stabilise a$, and so $[\inv g]_{\stabeq a} =
  [\unit]_{\stabeq a}$. Calculate:
  \[
  g \iact[\G/\A] a (f a)
  =
  g \iact[\G/\A] a [\unit]_{\stabeq a}
  =
  \bracks{\unit\mul\inv g}_{\stabeq{g \act a}}
  \explain={$g \in \stabilise a$}
  \bracks{\inv g}_{\stabeq a}
  =
  [\unit]_{\stabeq a}
  =
  f a
  \]
  Therefore $g \in \stabilise{f a}$, and so $\stabilise a \subset
  \stabilise{f a}$, and so $f$ is stable.

  By assumption, $f$ is equivariant, and by the Indexed Representation
  \theoremref{indexed-wlog-construction} we have an equivariant
  $\wlg(\canoniser,f) : \EFun\x\A (\G/\A)\x$. We will show that
  $\wlg(\canoniser,f) = \parent{\fun a\bracks{\canoniser a}_{\stabeq a}}$ and therefore
  $\canoniser$ is equivariant.

  Take any $a \in \carrier\A$ and calculate:
  \[
  \wlg(\canoniser, f)a = \inv*{\canoniser a}\act f(\canoniser a\act a)
  =
  \inv*{\canoniser a}\act \bracks\unit_{\stabeq{\canoniser a\act a}}
  =
  \bracks{\inv*{\inv*{\canoniser a}}}_{\stabeq{(\inv*{\canoniser a}) \act \canoniser a\act a}}
  =
  \bracks{\canoniser a}_{\stabeq a}
  \]
  Therefore $\wlg(\canoniser,f) = \bracks{\canoniser a}_{\stabeq a}$, as we wanted.
\end{proof}

We conclude this section by proving that every action has an
equivariant canoniser. Recall that a \emph{transversal} of an action
is a (dependent) function $t : (a : \carrier \A) \to \G\act\bracks a$
that picks a representative for the orbit of its input: $t a = t b$
iff $\G\act\bracks a = \G\act\bracks b$.

\begin{proof}[\theoremref{equivariant canoniser choice}]%
  Consider the family of orbits $\x : \A \types \OrbitOf\x$, and use
  the axiom of choice to choose a transversal: $t : \DFun\x\A
  \OrbitOf\x$.  For each $a \in \carrier\A$, the elements $a$ and $t\,a \in
  \OrbitOf a$ are symmetric, and so there is some symmetry $g \in
  \carrier\G$ such that $g \act a = t\,a$. Use the axiom of choice to
  choose a canoniser $\canoniser : \A \to \carrier\G$ such that
  $\canoniser a \act a = t\,a$. We will show this canoniser is
  equivariant.

  Take any element $a \in \carrier\A$ and symmetry $g \in \carrier\G$.
  Then $g \act a$ and $a$ are in the same orbit, and therefore $t(g \act a) = t\,a$.
  Calculate:
  \[
  \parent{\canoniser(g \act a) \mul g}\act a
  =
  \canoniser(g \act a) \act (g\act a)
  =
  t(g \act a)
  =
  t\,a
  =
  \canoniser a \act a
  \]
  Therefore, $\canoniser(g \act a) \mul g \stabeq a \canoniser a$, and so
  $\bracks{\canoniser(g \act a) \mul g}_{\stabeq a} = \bracks{\canoniser a}_{\stabeq a}$.
  Calculate:
  \[
  \bracks{\canoniser(g\act a)}_{\stabeq{g \act a}}
  =
  \bracks{\canoniser(g\act a)\mul g \mul \inv g}_{\stabeq{g \act a}}
  =
  g \act \bracks{\canoniser(g\act a)\mul g}_{\stabeq a}
  =
  g \act \bracks{\canoniser a}_{\stabeq a}
  \]
  and so $\canoniser$ is equivariant.
\end{proof}
